\subsection{Сложный Архимеда}
%%%%%%%%%%%%%%%%%%%%%%%%%%%%%%%%%%%%%%%%%%%%%%%%%%%%
%%%%%%%%%%%%%%%%%%%%%%%%%%%%%%%%%%%%%%%%%%%%%%%%%%%%
\z{В небольшом бассейне с площадью сечения $S$ плавает лодка, имеющая форму прямоугольного параллелепипеда с площадью нижней грани $S_\ell$. В лодке лежит камень массой $m$ и плотностью $\rho_k$. Увеличится или уменьшится уровень воды в бассейне, если камень выбросить на дно? На сколько?
}
%
{$\Delta h = \frac{m(\rho_k - \rho_v)}{\rho_v \rho_k S}.$}

%%%%%%%%%%%%%%%%%%%%%%%%%%%%%%%%%%%%%%%%%%%%%%%%%%%%
\z[{\begin{wrapfigure}[7]{r}{0.2\textwidth}%
		\vspace{-10pt}%
		\begin{flushright}%
 			\includegraphics[width=0.2\textwidth]{4.4.2}%
		\end{flushright}%
	\end{wrapfigure}%
	}]%
{Два кубика, связанные натянутой нитью, находятся в воде (см. рисунок). Верхний кубик со стороной $a = 10 \,\text{см}$ плавает, погрузившись в воду на три четверти своего объёма. Нижний кубик касается дна (вода под него подтекает). Сторона нижнего кубика равна $a/2$, а его плотность в 2 раза больше, чем у верхнего. Определите, при каких значениях плотности материала верхнего кубика возможно такое состояние системы. Плотность воды $\rho_0 = 1000 \,\text{кг/м}^3$.
\vspace{1em}
}
%
{$700 \,\text{кг/м}^3 < \rho < 750 \,\text{кг/м}^3$}

%%%%%%%%%%%%%%%%%%%%%%%%%%%%%%%%%%%%%%%%%%%%%%%%%%%%
\z[{\begin{wrapfigure}[5]{r}{0.2\textwidth}%
		\vspace{-30pt}%
		\begin{flushright}%
 			\includegraphics[width=0.2\textwidth]{4.4.3}%
		\end{flushright}%
	\end{wrapfigure}%
	}]%
{В цилиндрическом сосуде с водой плавает льдинка с привязанной к ней детской игрушкой. Силы натяжения всех нитей одинаковы и равны $T$. Определите, в какую сторону и на сколько изменится уровень воды в стакане после того, как лёд растает. Площадь дна сосуда $S$, плотность воды $\rho$.
\vspace{1em}
}
%
{$\Delta h = \frac{T}{\rho g h}$}

%%%%%%%%%%%%%%%%%%%%%%%%%%%%%%%%%%%%%%%%%%%%%%%%%%%%
\z[{\begin{wrapfigure}[9]{r}{0.2\textwidth}%
		\vspace{-15pt}%
		\begin{flushright}%
 			\includegraphics[width=0.2\textwidth]{4.4.4}%
		\end{flushright}%
	\end{wrapfigure}%
	}]%
{В цилиндрический сосуд с водой опущен металлический цилиндр, подвешенный на пружине жёсткостью $k = 10 \,\text{Н/м}$. Уровень воды в сосуде совпадает с положением верхней поверхности цилиндра (см. рисунок). После того как точку подвеса пружины подняли вверх на $y = 5 \,\text{см}$, удлинение пружины увеличилось на $x = 2 \,\text{см}$, а верхняя поверхность цилиндра оказалась на $h = 4 \,\text{см}$ выше уровня воды (нижняя поверхность цилиндра всё ещё в воде). Чему равна площадь поперечного сечения сосуда? 
\vspace{1em}
}
%
{$S = \frac{kx}{\rho g (x + h - y)} = 20 \,\text{см}^2$}

%%%%%%%%%%%%%%%%%%%%%%%%%%%%%%%%%%%%%%%%%%%%%%%%%%%%
\z[{\begin{wrapfigure}[10]{r}{0.2\textwidth}%
		\vspace{-0pt}%
		\begin{flushright}%
 			\includegraphics[width=0.2\textwidth]{4.4.5}%
		\end{flushright}%
	\end{wrapfigure}%
	}]%
{В герметичном сосуде сверху находится жидкость с плотностью $\rho_0 = 800 \,\text{кг/м}^3$, отделённая лёгким подвижным поршнем от газа (см. рисунок), находящегося внизу и имеющего давление $p = 20 \,\text{кПа}$. В поршне есть круглое отверстие, в которое вставлен цилиндрический поплавок, причём в жидкость поплавок погружён на длину $h$, а в газ — на длину $3h$. Площадь основания поплавка $S$. Поплавок может свободно скользить относительно поршня, а поршень -- относительно стенок сосуда. Жидкость нигде не подтекает. Какой должна быть плотность поплавка $\rho$, чтобы система могла оставаться в равновесии? 
\newpage
}
%
{$\rho = \frac{\rho_0}{4}$}

%%%%%%%%%%%%%%%%%%%%%%%%%%%%%%%%%%%%%%%%%%%%%%%%%%%%
\z[{\begin{wrapfigure}[7]{r}{0.25\textwidth}%
		\vspace{-0pt}%
		\begin{flushright}%
 			\includegraphics[width=0.25\textwidth]{4.4.6}%
		\end{flushright}%
	\end{wrapfigure}%
	}]%
{Два кубика, связанные натянутой нитью, находятся в воде. Верхний (со стороной $a = 1 \,\text{м}$) погружен в воду на три четверти объёма, а нижний — касается дна, и вода под него не подтекает. Сторона нижнего кубика равна $a/2$, а его плотность в 2 раза больше, чем у верхнего. Определите, при каких значениях плотности материала верхнего кубика возможно такое состояние системы. Толщина слоя воды $h = 2 \,\text{м}$. Внешнее давление $p_0 = 10^4 \,\text{Па}$.
}
%
{$100 \,\text{кг/м}^3 < \rho < 750 \,\text{кг/м}^3$}

%%%%%%%%%%%%%%%%%%%%%%%%%%%%%%%%%%%%%%%%%%%%%%%%%%%%
\z[{\begin{wrapfigure}[7]{r}{0.4\textwidth}%
		\vspace{-10pt}%
		\begin{flushright}%
 			\includegraphics[width=0.4\textwidth]{4.4.7}%
		\end{flushright}%
	\end{wrapfigure}%
	}]%
{На левую чашу находящихся в равновесии весов помещён сосуд, в который налито $V = 200 \,\text{мл}$ воды, а на правой — штатив с подвешенным на нити и погруженным в воду шариком. Объём шарика $V_\text{ш} = 20 \,\text{см}^3$. Какой груз и на какую чашу нужно положить, чтобы восстановить равновесие в случае замены воды в сосуде на равное по объёму количество керосина?
}
%
{$\text{груз массой } 48 \,\text{г на левую чашу}$}
